\chapter{Introduzione}
\label{ch:intro}

Con l’acronimo V2X si indica l’insieme delle comunicazioni tra un veicolo e ciò che lo circonda: altri veicoli (V2V), infrastruttura stradale (V2I) e servizi remoti. In termini pratici, queste comunicazioni permettono ai mezzi di condividere informazioni utili a guidare in modo più coordinato e sicuro, superando i limiti della sola percezione locale \cite{dressler2019cooperative,segata2016safe}.

Uno degli scenari più rappresentativi è il \emph{platooning}, cioè la marcia coordinata di più veicoli a distanza ridotta. In questa configurazione, un veicolo \emph{leader} guida il gruppo e i \emph{follower} lo seguono combinando i propri sensori con i dati ricevuti via rete. Se ben gestito, il platooning può migliorare capacità stradale, efficienza energetica e regolarità del traffico \cite{robinson2010operating,segata2016safe}.

Questi benefici dipendono però dalla qualità delle comunicazioni: ritardi o perdite di pacchetto possono influenzare direttamente stabilità e sicurezza del gruppo. Per questo, la letteratura recente considera rete, controllo e dinamica veicolare come parti di un unico problema da analizzare in modo integrato \cite{giordano2019joint,segata2023multi-technology}. In questa direzione si inserisce il presente lavoro, che viene inquadrato in modo più dettagliato nella sezione seguente.
